\section{Convergence analysis of the OSGS method}\label{app:osgs}

This appendix establishes, for the linearized OSGS variant, the stability
(\cref{oa:th:stability}) and convergence (\cref{oa:th:convergence}) results stated as
\cref{th:StabilityOSGS,th:ConvergenceOSGS} in \cref{sec:StabilityOSGS}. It follows the orthogonal-subscale template
of~\cite{Codina2008FiniteEA}, transported to the porosity-weighted operator, the Darcy
reaction term, and the artificial-compressibility perturbation. The proof is developed in
full; only the material shared with the ASGS analysis is not re-derived. Specifically, the
continuous problem, the Galerkin bilinear form and its coercivity identity, the stabilization
parameters, the mesh and interpolation-error setting, and the common standing assumptions are
stated once in the main text (\cref{sec:CommonSetting,sec:StabilizationDesign}) and, for the
ASGS route, in \cref{app:Continuity}; here they are only recalled and referenced, so that this
appendix develops the ingredients specific to the OSGS route---the orthogonal projection, the
weighted best-approximation and smoothing lemmas, the inf--sup argument, and the
Damk\"ohler-weighted error functional. Differences from the ASGS analysis of
\cref{app:Continuity} are flagged where they arise.

\subsection{The linearized problem and the OSGS method}
\label{oa:sec:setting}

\subsubsection{Continuous problem and Galerkin form}
\label{oa:sec:continuous}

The analyzed linearized problem, the Galerkin bilinear form $B(\ba;U,V)$
(\cref{eq:BilinearForm}) and the composite first-order quantity
$X(V) \coloneqq \ba\cdot\nabla\bv + \nabla q$ are those of the common setting
(\cref{sec:CommonSetting}) and of the ASGS analysis (\cref{sec:setting} of
\cref{app:Continuity}, where $X$ is denoted as in \cref{eq:XG}); we use them without
restatement. We recall only that, in the porous setting, the momentum residual carries the
weight $\alpha$, so the natural stabilized quantity is $\alpha X(V)$, and that
$(\nabla\bv, \alpha\nu\ViscProj\nabla\bu) = (\ViscProj\nabla\bv, \alpha\nu\ViscProj\nabla\bu)$
since $\ViscProj$ is a pointwise orthogonal projection and $\alpha\nu$ is scalar. We write
$(\cdot,\cdot)$ and $\norm{\cdot}$ for the $L^2(\Omega)$ inner product and norm (of scalars,
vectors or tensors), $\normK{\cdot} \coloneqq \norm{\cdot}_{L^2(K)}$,
$\norm{\cdot}_{m,K} \coloneqq \norm{\cdot}_{H^m(K)}$; the velocity and pressure spaces are
$V_0 \coloneqq H_0^1(\Omega)^d$ and $Q_0$ (the zero-mean subspace of $L^2(\Omega)$; for
$\varepsilon > 0$ the zero mean is automatic, since integrating the mass equation
$\varepsilon p + \nabla\cdot(\alpha\bu) = 0$ over $\Omega$ and using $\bu \in V_0$ gives
$\varepsilon\int_\Omega p = 0$), and $\mathcal{X}_0 \coloneqq V_0\times Q_0$.

Testing $B$ with $V = U$ and using $\nabla\cdot(\alpha\ba) = 0$ together with
$\bu|_{\partial\Omega} = \boldsymbol{0}$---so that the pressure--divergence pair cancels by
integration by parts and the convective term vanishes---gives the Galerkin coercivity
identity of the main text: for every $U \in \mathcal{X}_0$ with $\bu \in H^1_0(\Omega)^d$,
\begin{equation}
  \label{oa:eq:CoercivityIdentity}
  B(\ba; U, U)
  = 2\nu\,\bigl\lVert \alpha^{1/2}\ViscProj\nabla\bu \bigr\rVert^2
  + \norm{\sigma^{1/2}\bu}^2
  + \varepsilon\norm{p}^2 .
\end{equation}
This is \cref{eq:GalerkinCoercivityIdentity} of the main text, recalled here in the notation
of this appendix; the derivation is given there.

\subsubsection{Meshes, finite element spaces and interpolation}
\label{oa:sec:meshes}

The mesh nondegeneracy, the inverse estimate \cref{eq:InverseEstimateFiniteOrderNorm} and the
finite element spaces are the common ones of the analysis (\cref{sec:CommonSetting}), shared
with the ASGS setting of \cref{app:Continuity}; the OSGS route additionally uses the patch
structure and element{\hyp}size surrogates of \cref{H:patch} and the interpolation operators
recalled below.

We write $\mathcal{F} = \{\mathcal{T}_h\}_{h>0}$ for the (nondegenerate, shape-regular) mesh
family, $h_K \coloneqq \operatorname{diam}(K)$, $h \coloneqq \max_K h_K$; the continuous
degree-$k_u,k_p$ spaces $V_h \subset H^1(\Omega)^d$ and
$Q_h \subset H^1(\Omega)\cap C^0(\overline\Omega)$ (constants in $Q_h$, $k_p \le k_u+1$); and
the constrained subspaces $V_{h0} \coloneqq V_h \cap H^1_0(\Omega)^d$,
$Q_{h0} \coloneqq Q_h \cap Q_0$, $\Xhz \coloneqq V_{h0}\times Q_{h0}$.

\emph{Patch structure and element{\hyp}size surrogates} (\cref{H:patch}). We use the patch
quasi{\hyp}uniformity $\underline{C}\, h_{K} \le h_{K'} \le \overline{C}\, h_K$ for all
$K' \subset S_K$~\cite{AinsworthOden} and the continuous surrogates $\chi_1(\cdot,h),
\chi_2(\cdot,h)$ of \cref{H:patch}, with the element length taken as
$h_{K} \coloneqq \chi_1(\boldsymbol{x}_K, h)$ (as in~\cite[Lemma~1]{Codina2008FiniteEA}); this
last choice alters each stabilization parameter by a bounded factor ($\le \chi_0^2$) relative
to the raw diameter and is immaterial to every estimate, so we do not distinguish the two
notationally.

\emph{Interpolation operators.} For $v \in H^{r}(\Omega)$ with $r \ge 2$ (hence
$v \in C^0(\overline\Omega)$ for $d \le 3$) the nodal interpolant $\hat v_h$ satisfies
\begin{equation}
  \label{oa:eq:NodalInterp}
  \norm{v - \hat v_h}_{m,K} \le C_I\, h_K^{n-m}\norm{v}_{n,K},
  \qquad 0 \le m \le n \le \min(r,\, k+1),
\end{equation}
with $k$ the polynomial degree of the target space, and it preserves homogeneous Dirichlet
data (this is the estimate of the ASGS analysis, \cref{eq:interp} of \cref{app:Continuity},
in the norm form used below). For $v \in H^{r}(\Omega)$ with $r \ge 1$ the Scott--Zhang
quasi-interpolant~\cite{ScottZhang} satisfies the same estimate with $\norm{v}_{n,K}$
replaced by $\norm{v}_{H^n(S_K)}$ and also preserves homogeneous boundary values.

\subsubsection{Stabilization parameters}
\label{oa:sec:parameters}

The stabilization parameters $\tns$, $\tau_{1,K}$ and $\tau_{2,K}$ are those of the main text
(\cref{eq:TauNavierStokes,eq:Tau1Final,eq:Tau2Final}), shared with \cref{app:Continuity}, with
$\alpha_K \coloneqq \sup_K \alpha$, $|\ba|_{\infty,K} \coloneqq \norm{\ba}_{L^\infty(K)}$ and
design constants $c_1, c_2 > 0$ subject to \cref{H:design} of \cref{sec:CommonSetting}. Specific to the OSGS analysis are the
two dimensionless numbers and the $\boldsymbol{\tau}$-weighted broken products introduced now.
Expanding $\tau_2$,
\begin{equation}
  \label{oa:eq:Tau2Expanded}
  \tau_{2,K}
  = \frac{\nu}{\alpha_K}\Bigl( 1 + \frac{c_2}{c_1}\,\Reh \Bigr),
  \qquad
  \Reh \coloneqq \frac{|\ba|_{\infty,K}\, h_K}{\nu},
\end{equation}
so that $\tau_{2,K} \ge \nu/\alpha_K$ always. We will also use the elementwise mesh
Damk\"ohler number $\Dah \coloneqq \sigma h_K^2/(\alpha_K\nu)$. Weighted broken inner
products and norms are defined, for $i = 1,2$, by
\begin{equation}
  \label{oa:eq:TauInner}
  (f, g)_{\tau_i} \coloneqq \sum_K \tau_{i,K}\,(f, g)_{K},
  \qquad
  \norm{f}_{\tau_i} \coloneqq (f,f)_{\tau_i}^{1/2},
\end{equation}
for functions that need only be square integrable elementwise; note
$\norm{f}_{\tau_i} = \lVert \tau_i^{1/2} f\rVert_h$ in the broken-norm notation
of~the main text.

\subsubsection{Projections and the analyzed OSGS method}
\label{oa:sec:method}

Let $\Pone$ denote the $\tau_1$-weighted $L^2$ projection onto the
\emph{unconstrained} velocity space $V_h$ (acting componentwise), $\Ptwo$ the
$\tau_2$-weighted projection onto the unconstrained pressure space $Q_h$, and
$\Poneo \coloneqq I - \Pone$, $\Ptwoo \coloneqq I - \Ptwo$ the complementary
(fluctuation) operators. We will also need $\Ponez$, the $\tau_1$-weighted projection
onto the \emph{constrained} space $V_{h0}$, which enters only the stability proof
through the special test function, exactly as in~\cite{Codina2008FiniteEA}. Projecting onto the
unconstrained spaces in the method matches both the implementation described
in~the main text (footnote to the OSGS system there) and the discussion
in~\cite[Remark~1]{Codina2008FiniteEA}: projecting onto the constrained spaces would simplify
the analysis (assumption \cref{H:projection} would not be needed) but produces spurious numerical
boundary layers.

The \emph{analyzed} OSGS method reads: find $U_h \in \Xhz$ such that
\begin{equation}
  \label{oa:eq:Method}
  B_{\mathrm{osgs}}(\ba; U_h, V_h) = L(V_h) \coloneqq (\bff, \bv_h)
  \qquad \forall\, V_h \in \Xhz,
\end{equation}
with
\begin{equation}
  \label{oa:eq:Bosgs}
  B_{\mathrm{osgs}}(\ba; U, V)
  \coloneqq
  B(\ba; U, V) + S_1(U, V) + S_2(U, V),
\end{equation}
\begin{equation}
  \label{oa:eq:Sdefs}
  S_1(U, V) \coloneqq \bigl( \Poneo[\alpha X(U)],\, \alpha X(V) \bigr)_{\tau_1},
  \qquad
  S_2(U, V) \coloneqq \bigl( \Ptwoo[\nabla\cdot(\alpha\bu)],\, \nabla\cdot(\alpha\bv) \bigr)_{\tau_2}.
\end{equation}
Since $(\Poneo z, \Pone w)_{\tau_1} = 0$ for all $z, w$, both stabilization terms are
symmetric and positive semidefinite:
\begin{equation}
  \label{oa:eq:Ssym}
\begin{aligned}
  S_1(U, V) &= \bigl( \Poneo[\alpha X(U)],\, \Poneo[\alpha X(V)] \bigr)_{\tau_1},
  \qquad \\
  S_2(U, V) &= \bigl( \Ptwoo[\nabla\cdot(\alpha\bu)],\, \Ptwoo[\nabla\cdot(\alpha\bv)] \bigr)_{\tau_2}.
\end{aligned}
\end{equation}
In particular, with \cref{oa:eq:CoercivityIdentity},
\begin{equation}
  \label{oa:eq:Diagonal}
  B_{\mathrm{osgs}}(\ba; U_h, U_h)
  = 2\nu\,\bigl\lVert \alpha^{1/2}\ViscProj\nabla\bu_h \bigr\rVert^2
  + \norm{\sigma^{1/2}\bu_h}^2
  + \varepsilon\norm{p_h}^2
  + \norm{\Poneo \xi_h}_{\tau_1}^2
  + \norm{\Ptwoo \delta_h}_{\tau_2}^2,
\end{equation}
where here and below we abbreviate
\begin{equation}
  \label{oa:eq:XiDelta}
  \xi_h \coloneqq \alpha X(U_h) = \alpha\bigl(\ba\cdot\nabla\bu_h + \nabla p_h\bigr),
  \qquad
  \delta_h \coloneqq \nabla\cdot(\alpha\bu_h).
\end{equation}

The next elementary observation locates the analyzed method precisely relative to the
implemented one, in which the orthogonal projection is applied to the \emph{full}
residual and the test slot is the full adjoint.

\begin{lemma}[Exact annihilations]
\label{oa:lem:annihilation}
Let $\sigma$ and $\varepsilon$ be constant. Then, for all $U_h, V_h \in \Xhz$:
\begin{enumerate}[label=\textup{(\roman*)}, itemsep=0.2ex]
  \item $\Poneo(\sigma\bu_h) = \boldsymbol{0}$ and $\Ptwoo(\varepsilon p_h) = 0$
        (reactive and compressibility parts of the residual have no fluctuation);
  \item $\bigl( \Poneo z,\, \sigma\bv_h \bigr)_{\tau_1} = 0$ and
        $\bigl( \Ptwoo w,\, \varepsilon q_h \bigr)_{\tau_2} = 0$ for any $z, w$
        (reactive and compressibility parts of the adjoint test slot are invisible to
        the fluctuations).
\end{enumerate}
Consequently, the formulation obtained from the implemented OSGS method by projecting
the full residual and testing with the full adjoint differs from
\cref{oa:eq:Method}--\cref{oa:eq:Sdefs} only through (a) the fluctuations of the viscous
residual term $2\nabla\cdot(\alpha\nu\ViscProj\nabla\bu_h)$ and of the force $\bff$,
and (b) the viscous part of the adjoint test slot---exactly the Method~I simplification
of~\cite{Codina2008FiniteEA}.
\end{lemma}

\begin{proof}
For (i), $\sigma\bu_h \in V_h$ and $\varepsilon p_h \in Q_h$, and the fluctuation of
any element of the projection space vanishes. For (ii), $\sigma\bv_h \in V_h$ and
$\varepsilon q_h \in Q_h$, and fluctuations are $\tau_i$-orthogonal to the projection
spaces. The final statement follows by expanding the implemented stabilization
term by linearity.
\end{proof}

\begin{remark}[Analyzed versus implemented formulation]
\label{oa:rem:analyzed}
For later reference we list the exact distance between \cref{oa:eq:Method} and the OSGS
method as implemented in~the main text. (i)~\emph{Weighted versus plain projection}:
the analysis uses the $\boldsymbol{\tau}$-weighted projections of
\cref{oa:eq:TauInner}; the implementation uses the plain $L^2$ projection. The two
coincide when $\tau_i$ is globally constant; otherwise the weighted product is the one
for which the best-approximation property of \cref{oa:lem:bestapprox} holds, which is
essential to the consistency estimate (cf.~\cite[Remark~4]{Codina2008FiniteEA}).
(ii)~\emph{First-order truncation}: residual and test slots are truncated to their
first-order parts; by \cref{oa:lem:annihilation} the only genuinely neglected pieces are
the viscous fluctuation, the force fluctuation, and the viscous test term. For $k_u = 1$
and elementwise-constant $\alpha\nu$ the viscous residual vanishes on element interiors;
with variable $\alpha$ it contributes $2\nu(\ViscProj\nabla\bu_h)\nabla\alpha$, which is
small under the resolution condition of \cref{H:porosity}. (iii)~\emph{Converged versus lagged
projection}: the analysis takes the projection equation solved exactly (the coupled
system), whereas the implementation lags the projection along the nonlinear iteration.
(iv)~\emph{Unconstrained projection space}: analyzed and implemented formulations
agree (both project onto $V_h\times Q_h$ without Dirichlet constraints); note that the
display defining the OSGS projection in~the main text shows the projection onto
$\Xhz$, while the accompanying footnote states the implemented choice, which is the one
analyzed here and the one recommended by~\cite[Remark~1]{Codina2008FiniteEA}.
(v)~\emph{Elementwise-constant versus pointwise parameters}: as for the ASGS analysis
of~the main text, $\tau_{i}$ is elementwise constant in the analysis and pointwise
variable in the computations.
\end{remark}

\subsection{Standing assumptions and notation}
\label{oa:sec:hypotheses}

Throughout, $C$ denotes a generic positive constant, independent of $h$, $\nu$,
$\sigma$, $\varepsilon$ and of $\alpha$ and $\ba$ except through the constants named
below; its value may change between appearances. The analysis is carried out under the
standing assumptions collected in \cref{sec:CommonSetting}: the assumptions
\cref{H:mesh}--\cref{H:regularity} common to both methods, and the OSGS{\hyp}specific
assumptions \cref{H:design}--\cref{H:advectionsmooth}---namely the design constants
\cref{H:design} (with the compressibility bound $\varepsilon\tau_{2,K}\le C_2$), the projection
stability \cref{H:projection} (\cref{eq:ProjectionStability}), and the technical strengthenings
\cref{H:patch} (patch structure and element{\hyp}size surrogates), \cref{H:porositysmooth}
(porosity smoothness) and \cref{H:advectionsmooth} (advection smoothness) of the common mesh,
porosity and advection assumptions. The strengthenings \cref{H:porositysmooth,H:advectionsmooth}
enter only the consistency estimate. The ASGS analysis of \cref{app:Continuity} instead requires
the interior{\hyp}face jump condition \cref{H:jump}; the two hypothesis sets are therefore not
nested (see \cref{oa:rem:dispensed}). We use the abbreviations $\Ponez z_h$ and $\Poneo z_h$ of
\cref{oa:sec:method} for the constrained projection and the fluctuation appearing in
\cref{H:projection}, and, for \cref{H:porositysmooth}, $C_{\alpha,m}$ ($m\ge2$) as defined there,
setting $C_{\alpha,1}\coloneqq 1+\norm{\nabla\alpha}_{L^\infty(\Omega)}/\alpha_0$ for the
first{\hyp}order case.

\subsection{Technical results}
\label{oa:sec:technical}

\subsubsection{Parameter inequalities}

Items (P1)--(P3) below are elementary parameter bounds shared with the ASGS analysis: they are
recalled from \cref{lem:parameters} of \cref{app:Continuity}---in the notation of this appendix,
through $\varphi_1=\alpha_K\tns^{-1}$---and are not re-derived. Items (P4)--(P6) and the
inverse{\hyp}estimate inequalities (K1)--(K3) are specific to the OSGS route (they feed the
special test function of \cref{oa:th:stability} and the Damk\"ohler-weighted conversions of the
convergence proof).

\begin{lemma}[Parameter inequalities]
\label{oa:lem:parameters}
Under \cref{H:design}, the parameters \cref{eq:Tau1Final,eq:Tau2Final} satisfy, on every element $K$
(element subscripts omitted inside the formulas):
\begin{align}
  &\sigma\tau_1 \le 1,
  \qquad
  \alpha_K \tau_1 \le \tns,
  \qquad
  \varepsilon \le C_2\, \tau_2^{-1},
  \tag{P1}\label{oa:eq:P1}\\
  &\tau_1\,\tau_2 \le \frac{h_K^2}{c_1\,\alpha_K^2},
  \tag{P2}\label{oa:eq:P2}\\
  &h_K^2\,\tns^{-1} = c_1\,\alpha_K\,\tau_2,
  \qquad\text{hence}\qquad
  \bigl(\alpha_K\,\tns^{-1}\bigr)^{1/2} = \sqrt{c_1}\,\frac{\alpha_K}{h_K}\,\tau_2^{1/2},
  \tag{P3}\label{oa:eq:P3}\\
  &\nu \le \alpha_K\,\tau_2,
  \qquad
  \tau_1^{1/2}\,|\ba|_{\infty,K} \le \frac{\sqrt{c_1}}{c_2}\,\tau_2^{1/2},
  \tag{P4}\label{oa:eq:P4}\\
  &\alpha_K^2\,\tau_1\,\tau_2 \;\ge\; \frac{h_K^2}{c_1 + \Dah},
  \qquad\text{hence}\qquad
  \tau_2^{-1/2} \le \bigl(c_1 + \Dah\bigr)^{1/2}\,\frac{\alpha_K}{h_K}\,\tau_1^{1/2},
  \tag{P5}\label{oa:eq:P5}\\
  &\sigma^{1/2} \;\le\; \Dah^{1/2}\,\frac{\alpha_K}{h_K}\,\tau_2^{1/2}.
  \tag{P6}\label{oa:eq:P6}
\end{align}
Moreover, with $\gamma$ as in \cref{H:design},
\begin{align}
  &C_{\mathrm{inv}}\,\frac{(\alpha_K\nu)^{1/2}}{h_K} \le \frac{1}{\gamma}\,\tau_1^{-1/2},
  \tag{K1}\label{oa:eq:K1}\\
  &C_{\mathrm{inv}}\,\frac{\alpha_K\,|\ba|_{\infty,K}}{h_K} \le \frac{1}{\gamma}\,\tau_1^{-1},
  \tag{K2}\label{oa:eq:K2}\\
  &C_{\mathrm{inv}}\,\frac{\alpha_K\,(\tau_1\tau_2)^{1/2}}{h_K} \le \frac{1}{\gamma}.
  \tag{K3}\label{oa:eq:K3}
\end{align}
\end{lemma}

\begin{proof}
The first two inequalities of \cref{oa:eq:P1} ($\sigma\tau_1\le1$ and $\alpha_K\tau_1\le\tns$)
and the identity \cref{oa:eq:P3} are established in \cref{lem:parameters} of
\cref{app:Continuity} (with $\varphi_1=\alpha_K\tns^{-1}$) and are not re-derived here; the
third inequality of \cref{oa:eq:P1}, $\varepsilon\le C_2\tau_2^{-1}$, is \cref{H:design}, and
\cref{oa:eq:P2} follows on combining $\alpha_K\tau_1\le\tns$ with the definition of $\tau_2$. We
prove the OSGS{\hyp}specific items. \cref{oa:eq:P4}: the first inequality is
\cref{oa:eq:Tau2Expanded}; for the second, $\tau_1 \le \tns/\alpha_K \le
h_K/(c_2\alpha_K|\ba|_{\infty,K})$ gives
$\tau_1|\ba|_{\infty,K}^2 \le |\ba|_{\infty,K}h_K/(c_2\alpha_K)$, while
\cref{oa:eq:Tau2Expanded} gives
$\tau_2 \ge (c_2/c_1)\,|\ba|_{\infty,K}h_K/\alpha_K$; divide. \cref{oa:eq:P5}: using
$\tns \le h_K^2/(c_1\nu)$,
\[
\begin{aligned}
  \alpha_K^2\tau_1\tau_2
  &= \frac{\alpha_K^2}{\alpha_K\tns^{-1}+\sigma}\cdot\frac{h_K^2}{c_1\alpha_K\tns}
  = \frac{h_K^2}{c_1\bigl(1 + \sigma\tns/\alpha_K\bigr)} \\
  &\ \ge\ \frac{h_K^2}{c_1\bigl(1 + \sigma h_K^2/(c_1\alpha_K\nu)\bigr)}
  \ =\ \frac{h_K^2}{c_1 + \Dah}.
\end{aligned}
\]
\cref{oa:eq:P6}: $\sigma^{1/2} = \Dah^{1/2}(\alpha_K\nu)^{1/2}/h_K \le
\Dah^{1/2}\alpha_K\tau_2^{1/2}/h_K$ by the first inequality of \cref{oa:eq:P4}.
\cref{oa:eq:K1}: $\tau_1^{-1} \ge \alpha_K\tns^{-1} \ge c_1\alpha_K\nu/h_K^2$, so
$C_{\mathrm{inv}}(\alpha_K\nu)^{1/2}/h_K \le
(C_{\mathrm{inv}}/\sqrt{c_1})\tau_1^{-1/2} \le \gamma^{-1}\tau_1^{-1/2}$ by \cref{H:design}.
\cref{oa:eq:K2}: $\tau_1^{-1} \ge \alpha_K\tns^{-1} \ge c_2\alpha_K|\ba|_{\infty,K}/h_K$
and $c_2 \ge \gamma C_{\mathrm{inv}}$. \cref{oa:eq:K3}: combine \cref{oa:eq:P2} with
$c_1 \ge \gamma^2C_{\mathrm{inv}}^2$.
\end{proof}

\subsubsection{Patch equivalence and the smoothing operation}

\begin{lemma}[Patch equivalence]
\label{oa:lem:patch}
Under \cref{H:mesh}, \cref{H:porosity}, \cref{H:advection} and \cref{H:patch} there are constants
$0 < \underline{c} \le \overline{c}$, independent of $h$, such that for $i = 1,2$,
\begin{equation}
  \label{oa:eq:PatchEquivalence}
  \underline{c}\;\tau_{i,K} \le \tau_{i,K'} \le \overline{c}\;\tau_{i,K},
  \qquad
  \underline{c}\;\alpha_{K} \le \alpha_{K'} \le \overline{c}\;\alpha_{K},
  \qquad \forall\, K' \subset S_K,\ \forall\, K.
\end{equation}
\end{lemma}

\begin{proof}
Patch quasi-uniformity gives $h_{K'} \eqsim h_K$. Since $\alpha$ is Lipschitz and
bounded below by $\alpha_0 > 0$, and $\operatorname{diam}(S_K) \le C h_K$, one has
$|\alpha_{K'} - \alpha_K| \le C h_K\norm{\nabla\alpha}_{L^\infty(S_K)} \le
C\,C_{\nabla\alpha}\,\max_{K''\subset S_K}\alpha_{K''}$, from which the uniform equivalence of
the $\alpha_{K}$ over the patch follows for $C\,C_{\nabla\alpha}$ bounded (and, in general, from
$\alpha_{K'}/\alpha_K \le \alpha_\infty/\alpha_0$; the sharper route through $C_{\nabla\alpha}$
keeps the constant independent of $\alpha_\infty/\alpha_0$). Likewise
$\ba \in C^0(\overline\Omega)$ implies $|\ba|_{\infty,K'} \le |\ba|_{\infty,K} +
\omega_{\ba}(Ch_K)$, where $\omega_{\ba}$ is the modulus of continuity, and the
resulting perturbations of $\tns^{-1}$, $\tau_1^{-1}$ and $\tau_2$ are controlled as
in the proof of \cref{oa:lem:smoothing} below. The claim follows from the explicit
formulas \cref{eq:Tau1Final,eq:Tau2Final}.
\end{proof}

Following~\cite{Codina2008FiniteEA}, for $\tau \in \{\tau_1, \tau_2\}$ and a finite element
function $v_h$ (scalar or vector) we define the discontinuous and the smoothed
products
\begin{equation}
  \label{oa:eq:SmulDef}
  (\tau \pmul v_h)\big|_K \coloneqq \tau_{K}\, v_h\big|_K,
  \qquad
  (\tau \smul v_h)\big|_K \coloneqq \sum_{a \in N_K} N_a(\boldsymbol{x})\big|_K\,
  \bar\tau_{a}\, v_a,
\end{equation}
where $N_a$ are the nodal shape functions, $v_a$ the nodal values of $v_h$, and the
nodal parameters $\bar\tau_{i,a}$ are obtained by evaluating the formulas
\cref{eq:Tau1Final,eq:Tau2Final} with $h_K$ replaced by $\chi_1(\boldsymbol{x}_a, h)$,
$\alpha_K$ by $\alpha(\boldsymbol{x}_a)$ and $|\ba|_{\infty,K}$ by
$|\ba(\boldsymbol{x}_a)|$. Thus $\tau \smul v_h$ is a continuous finite element
function, whereas $\tau \pmul v_h$ is not.

\begin{lemma}[Smoothing; adaptation of~{\cite[Lemma~2]{Codina2008FiniteEA}}]
\label{oa:lem:smoothing}
Under \cref{H:mesh}, \cref{H:porosity}, \cref{H:advection} and \cref{H:patch} there is a function $\psi(h) \to 0$
as $h \to 0$ such that, for $\tau \in \{\tau_1, \tau_2\}$ and any finite element
function $v_h$,
\begin{equation}
  \label{oa:eq:Smoothing}
  \norm{\tau\pmul v_h - \tau\smul v_h}_{K} \le \tau_K\,\psi(h)\,\normK{v_h},
  \qquad\text{hence}\qquad
  \norm{\tau\smul v_h}_{K} \le (1+\psi(h))\,\tau_K\,\normK{v_h}.
\end{equation}
If, in addition, $\ba$ is Lipschitz on $\overline\Omega$, then $\psi(h) = O(h)$.
\end{lemma}

\begin{proof}
As in~\cite[Lemma~2]{Codina2008FiniteEA}, the equivalence
$C_1 h_K^{d/2}\norm{w_h}_{L^\infty(K)} \le \normK{w_h} \le
C_2 h_K^{d/2}\norm{w_h}_{L^\infty(K)}$ for finite element functions on nondegenerate
meshes reduces \cref{oa:eq:Smoothing} to the statement
$\max_{a\in N_K} |\bar\tau_a/\tau_K - 1| \le C\,\psi(h)$ with $\psi(h)\to 0$
uniformly in $K$. Write $\Delta(\cdot)$ for the difference between the nodal and the
elemental value of a quantity. For $\tau_1$,
$\tau_{1}|\Delta(\tau_1^{-1})| \le
\tau_1\bigl[\,|\Delta\alpha|\,\tns^{-1}
+ \alpha_\infty\, c_2\,|\Delta|\ba||/h_K
+ \alpha_K\, c_1\nu\,|\Delta(h^{-2})| + \alpha_K c_2 |\ba|_{\infty,K}|\Delta(h^{-1})|\,\bigr]$.
Using $\tau_1 \le \tns/\alpha_0$ and $\tau_1 \le h_K^2/(c_1\alpha_0\nu)$,
the first term is bounded by $\omega_\alpha(Ch)/\alpha_0$, the second by
$\alpha_\infty\,(c_2/(c_1\alpha_0\nu))\,h\,\omega_{\ba}(Ch)$, and the last two by
$C\,\omega_{\chi}(h)$, where $\omega_\alpha, \omega_{\ba}$ are the moduli of
continuity of $\alpha$ and $\ba$ and $\omega_\chi(h)\to 0$ quantifies the relative
oscillation of $\chi_1(\cdot,h)$ over an element, which vanishes as $h \to 0$ by the
construction of~\cite[Lemma~1]{Codina2008FiniteEA}. All right-hand sides tend to zero
uniformly; here the strict positivity $\alpha_0 > 0$ and $\nu > 0$ are used, and for
Lipschitz data all moduli are $O(h)$. Then
$|\bar\tau_{1,a}/\tau_{1,K} - 1| = \bar\tau_{1,a}\,|\Delta(\tau_1^{-1})|
\le (1+o(1))\,\tau_{1,K}|\Delta(\tau_1^{-1})|$, which proves the claim for $\tau_1$;
the same computation applied to the explicit form \cref{oa:eq:Tau2Expanded} of $\tau_2$
(a continuous, positive expression in $\alpha$, $|\ba|$ and $h$, bounded away from
zero by $\nu/\alpha_\infty$) proves it for $\tau_2$.
\end{proof}

\subsubsection{Best approximation by the weighted projections}

\begin{lemma}[Best approximation and weighted residual approximation]
\label{oa:lem:bestapprox}
Let $i \in \{1,2\}$ and let $\Pi_i$ be as in \cref{oa:sec:method}. Then, for any
elementwise square-integrable $z$ and any finite element function $z_h$ in the
corresponding projection space,
\begin{equation}
  \label{oa:eq:BestApprox}
  \norm{\Pi_i^{\perp} z}_{\tau_i} \le \norm{z - z_h}_{\tau_i}.
\end{equation}
Moreover, let $z \in H^{r}(\Omega)^s$ with $1 \le r \le k_u+1$ (for $i = 1$) or
$1 \le r \le k_p + 1$ (for $i = 2$), and let $\alpha$ satisfy \cref{H:porositysmooth} with
$m \ge r$ wherever $r \ge 2$. Then
\begin{equation}
  \label{oa:eq:WeightedApprox}
  \norm{\Pi_i^{\perp}(\alpha z)}_{\tau_i}
  \;\le\;
  C\, C_{\alpha,r}
  \Bigl( \sum_K \tau_{i,K}\,\alpha_K^2\, h_K^{2r}\,\norm{z}_{H^{r}(K)}^2 \Bigr)^{1/2},
\end{equation}
with $C$ depending only on the interpolation constant, the patch-overlap constant and
the constants of \cref{oa:lem:patch}, and $C_{\alpha,r}$ as in \cref{H:porositysmooth}
(with $C_{\alpha,1} \coloneqq 1 + \norm{\nabla\alpha}_{L^\infty(\Omega)}/\alpha_0$).
\end{lemma}

\begin{proof}
\cref{oa:eq:BestApprox} is the minimizing property of orthogonal projections in their own
inner product, cf.~\cite[Remark~4]{Codina2008FiniteEA}. For \cref{oa:eq:WeightedApprox}, take
$z_h$ to be the Scott--Zhang interpolant of $\alpha z$ in the projection space, so that
$\norm{\alpha z - z_h}_{K} \le C h_K^{r} \norm{\alpha z}_{H^{r}(S_K)}$
(here $r \le k+1$ for the relevant degree $k$, by \cref{H:spaces} and the stated restrictions on
$r$; in particular the case $i=1$, $z = \nabla p$, $r = k_p$ uses $k_p \le k_u+1$).
By the Leibniz rule and \cref{H:porositysmooth},
$\norm{\alpha z}_{H^{r}(S_K)}
\le \bigl( \alpha_{\infty,S_K} + \sum_{j=1}^{r}\binom{r}{j}\norm{D^j\alpha}_{L^\infty}\bigr)
\norm{z}_{H^{r}(S_K)}
\le C\,C_{\alpha,r}\,\alpha_{\infty,S_K}\norm{z}_{H^{r}(S_K)}$,
using $\norm{D^j\alpha}_{L^\infty} \le (C_{\alpha,r}-1)\alpha_0 \le
(C_{\alpha,r}-1)\alpha_{\infty,S_K}$. Insert this into \cref{oa:eq:BestApprox}, replace
$\alpha_{\infty,S_K}$ and $\tau_{i,S_K}$-values by their values on $K$ using
\cref{oa:lem:patch}, and re-sum over patches using the bounded overlap implied by
nondegeneracy, absorbing $\norm{z}_{H^r(S_K)}$ into element contributions.
\end{proof}

\subsubsection{The OSGS triple norm}
\label{oa:sec:norm}

The stability and convergence analysis is carried out in the mesh-dependent OSGS norm of the
main text (\cref{eq:TripleNormOSGS}), recalled here for convenience,
\begin{equation}
  \label{oa:eq:TripleNorm}
  \tnorm{V}_{\mathrm O}
  \coloneqq
  \Bigl(
  2\nu\,\bigl\lVert \alpha^{1/2}\ViscProj\nabla\bv \bigr\rVert^2
  + \norm{\sigma^{1/2}\bv}^2
  + \varepsilon\norm{q}^2
  + \norm{\alpha X(V)}_{\tau_1}^2
  + \norm{\nabla\cdot(\alpha\bv)}_{\tau_2}^2
  \Bigr)^{1/2}.
\end{equation}
Throughout this appendix we keep the method subscripts of the main text:
$\tnorm{\cdot}_{\mathrm O}$ denotes this OSGS norm (\cref{eq:TripleNormOSGS}), while
$\tnorm{\cdot}_{\mathrm S}$ denotes the ASGS norm of \cref{lemma:Stability} of the main text,
which controls only the damped reaction $\sigmatilde$ (\cref{oa:lem:normcomparison}). The bare
triple bar is not used.

\begin{lemma}[Definiteness]
\label{oa:lem:definiteness}
Under \cref{H:mesh}--\cref{H:porosity}, and with the pressure normalized to zero mean when
$\varepsilon = 0$, \cref{oa:eq:TripleNorm} is a norm on $\Xhz$; in particular $\tnorm{V}_{\mathrm{O}} = 0$
implies $V = 0$.
\end{lemma}

\begin{proof}
Homogeneity, symmetry and the triangle inequality are immediate; only definiteness needs an
argument, and it is here that the full Dirichlet data enter. Let $\tnorm{V}_{\mathrm{O}} = 0$ for
$V = [\bv; q] \in \Xhz$. Then $\ViscProj\nabla\bv \equiv 0$, i.e.\ the symmetric gradient of
$\bv$ is spherical almost everywhere; since $\bv \in H^1_0(\Omega)^d$, the deviatoric
(conformal) Korn inequality on the bounded domain $\Omega$---whose conformal--Killing kernel
is excluded by the homogeneous boundary data---gives $\norm{\nabla\bv} \le
C\,\bigl\lVert\ViscProj\nabla\bv\bigr\rVert = 0$, so $\bv \equiv 0$. With $\bv = 0$, the term
$\norm{\alpha X(V)}_{\tau_1} = 0$ reduces to $\norm{\alpha\nabla q}_{\tau_1} = 0$; as
$\alpha \ge \alpha_0 > 0$ and $\tau_1 > 0$, this forces $\nabla q \equiv 0$, so $q$ is
constant, and the normalization ($\varepsilon\norm{q} = 0$ when $\varepsilon > 0$, or the
zero-mean constraint when $\varepsilon = 0$) forces $q \equiv 0$. Hence $\tnorm{\cdot}_{\mathrm{O}}$ is
definite. (Absent full Dirichlet data or pressure normalization it is only a seminorm.)
\end{proof}

\begin{lemma}[Comparison with the ASGS norm]
\label{oa:lem:normcomparison}
Let $\tnorm{\cdot}_{\mathrm{S}}$ denote the triple norm of the ASGS stability lemma
of~the main text,
\[
\begin{aligned}
  \tnorm{V}_{\mathrm{S}}^2
  &= \nu\,\bigl\lVert \alpha^{1/2}\ViscProj\nabla\bv \bigr\rVert^2
  + \bigl\lVert \widetilde{\sigma}_\alpha^{1/2}\bv \bigr\rVert_h^2
  + \varepsilon\norm{q}^2
  + \norm{\alpha X(V)}_{\tau_1}^2
  + \norm{\nabla\cdot(\alpha\bv)}_{\tau_2}^2,
  \qquad \\
  \widetilde{\sigma}_\alpha
  &= \frac{\tns^{-1}\sigma}{\tns^{-1} + \sigma/\alpha_K}.
\end{aligned}
\]
Then $\tnorm{V}_{\mathrm{S}} \le \tnorm{V}_{\mathrm{O}}$ for all $V$.
\end{lemma}

\begin{proof}
Term by term: $\nu \le 2\nu$; on each element
$\widetilde{\sigma}_\alpha = \sigma\,\tns^{-1}/(\tns^{-1}+\sigma/\alpha_K) \le \sigma$;
the remaining three terms coincide.
\end{proof}

\subsection{Stability}
\label{oa:sec:stability}

\begin{theorem}[Inf--sup stability]
\label{oa:th:stability}
Under \cref{H:mesh}--\cref{H:spaces}, \cref{H:design}--\cref{H:patch} there exist $h_0 > 0$ and $\beta_{\mathrm{osgs}} > 0$,
depending only on $\beta_0$, $C_2$, $C_{\nabla\alpha}/C_{\mathrm{inv}}$ and the constants of
\cref{oa:sec:technical}, such that for all $h \le h_0$,
\begin{equation}
  \label{oa:eq:InfSup}
  \inf_{U_h \in \Xhz\setminus\{0\}}\;
  \sup_{V_h \in \Xhz\setminus\{0\}}\;
  \frac{B_{\mathrm{osgs}}(\ba; U_h, V_h)}{\tnorm{U_h}_{\mathrm{O}}\,\tnorm{V_h}_{\mathrm{O}}}
  \;\ge\; \beta_{\mathrm{osgs}} .
\end{equation}
An admissible choice of the constants in \cref{H:design} is
$\gamma_0 = 10\,(1 + M_2)$ with $M_2 \coloneqq 1 + C_{\nabla\alpha}/C_{\mathrm{inv}}$,
together with $h_0$ small enough that
$\psi(h_0) \le \min\{1/8,\ \beta_0^2/16\}$, where $\psi$ is the function of
\cref{oa:lem:smoothing}. These are \emph{sufficient} conditions, and $\gamma_0$ is not claimed
sharp: with $\gamma_0 \ge 20$ they demand $c_1 \gtrsim 400\,C_{\mathrm{inv}}^2$ and
$c_2 \gtrsim 20\,C_{\mathrm{inv}}$, well above the values $c_1 = 4k^4$ (or $16k^4$),
$c_2 = 2k^2$ used in the experiments of~the main text, which are not shown to satisfy
them. The result is thus a \emph{conditional} inf--sup stability theorem, not a calibration
or validation of the stabilization constants employed in the computations; the
threshold~$h_0$ likewise depends on the coefficient data through $\psi$
(\cref{oa:lem:smoothing}) and is not uniform in $\alpha_0$, $\nu$ or the moduli of continuity
of $\alpha$ and $\ba$.
\end{theorem}

\begin{proof}
Fix $U_h = [\bu_h; p_h] \in \Xhz$ and recall the abbreviations \cref{oa:eq:XiDelta}.
Write $\psi \coloneqq \psi(h)$ and set
\[
  D_\nu \coloneqq 2\nu\,\bigl\lVert \alpha^{1/2}\ViscProj\nabla\bu_h \bigr\rVert^2,
  \quad
  D_\sigma \coloneqq \norm{\sigma^{1/2}\bu_h}^2,
  \quad
  D_\varepsilon \coloneqq \varepsilon\norm{p_h}^2,
\]
\[
  O_1 \coloneqq \norm{\Poneo\xi_h}_{\tau_1}^2,
  \quad
  O_2 \coloneqq \norm{\Ptwoo\delta_h}_{\tau_2}^2,
  \quad
  P_1 \coloneqq \norm{\Ponez\xi_h}_{\tau_1}^2,
  \quad
  P_2 \coloneqq \norm{\Ptwo\delta_h}_{\tau_2}^2 .
\]

\medskip\noindent\emph{Step 1: the special test function.}
Define
\[
  \bw_h \coloneqq \tau_1 \smul \Ponez(\xi_h) \;\in\; V_{h0},
  \qquad
  r_h \coloneqq \tau_2 \smul \Ptwo(\delta_h) - m_h \;\in\; Q_{h0},
\]
where $m_h$ is the mean of $\tau_2\smul\Ptwo(\delta_h)$ over $\Omega$ whenever
$Q_{h0}$ consists of zero-mean functions, and $m_h \coloneqq 0$ otherwise. Note that
$\bw_h$ is an admissible velocity test function ($\Ponez\xi_h \in V_{h0}$ vanishes on
$\partial\Omega$, and the smoothed product preserves nodal zeros), and $r_h$ an
admissible pressure test function since constants belong to $Q_h$ by \cref{H:spaces}. Set
$V^0 \coloneqq [\bw_h; r_h]$ and take as test function $W \coloneqq U_h + V^0$.
By \cref{oa:eq:Diagonal},
\begin{equation}
  \label{oa:eq:StepDiag}
  B_{\mathrm{osgs}}(\ba; U_h, U_h) = D_\nu + D_\sigma + D_\varepsilon + O_1 + O_2 ,
\end{equation}
and it remains to bound $B_{\mathrm{osgs}}(\ba; U_h, V^0)$ from below. Expanding
\cref{eq:BilinearForm} and \cref{oa:eq:Sdefs},
\begin{align}
  B_{\mathrm{osgs}}(\ba; U_h, V^0)
  ={}& 2\nu\,(\alpha\,\ViscProj\nabla\bu_h, \ViscProj\nabla\bw_h)
  \tag{T1}\label{oa:eq:T1}\\
  &+ (\bw_h, \xi_h)
  \tag{T2}\label{oa:eq:T2}\\
  &+ \sigma(\bu_h, \bw_h)
  \tag{T3}\label{oa:eq:T3}\\
  &+ \varepsilon(p_h, r_h)
  \tag{T4}\label{oa:eq:T4}\\
  &+ (r_h, \delta_h)
  \tag{T5}\label{oa:eq:T5}\\
  &+ \bigl(\Poneo\xi_h,\ \alpha(\ba\cdot\nabla\bw_h + \nabla r_h)\bigr)_{\tau_1}
  \tag{T6}\label{oa:eq:T6}\\
  &+ \bigl(\Ptwoo\delta_h,\ \nabla\cdot(\alpha\bw_h)\bigr)_{\tau_2},
  \tag{T7}\label{oa:eq:T7}
\end{align}
where in \cref{oa:eq:T2} we used
$(\bw_h, \alpha\ba\cdot\nabla\bu_h) + (\bw_h, \alpha\nabla p_h) = (\bw_h, \xi_h)$.
The constant shift $m_h$ is invisible throughout: it does not affect $\nabla r_h$ in
\cref{oa:eq:T6}; in \cref{oa:eq:T5} it contributes $-m_h\int_\Omega \delta_h = 0$ because
$\int_\Omega \nabla\cdot(\alpha\bu_h) = 0$ by the divergence theorem and
$\bu_h \in V_{h0}$; and in \cref{oa:eq:T4} it contributes
$-\varepsilon\, m_h \int_\Omega p_h = 0$, since either $m_h = 0$ or $p_h$ has zero
mean. We now estimate \cref{oa:eq:T1}--\cref{oa:eq:T7} term by term. Throughout,
\cref{oa:eq:Smoothing} is used in the form
$\normK{\bw_h} \le (1+\psi)\,\tau_{1,K}\normK{\Ponez\xi_h}$ and
$\normK{r_h + m_h} \le (1+\psi)\,\tau_{2,K}\normK{\Ptwo\delta_h}$ (recall also that
subtracting the mean does not increase the $L^2$ norm, so the same bound holds
globally for $r_h$).

\medskip\noindent\emph{Step 2: the two productive terms \cref{oa:eq:T2} and \cref{oa:eq:T5}.}
Splitting $\tau_1\smul(\cdot) = \tau_1\pmul(\cdot) + (\tau_1\smul - \tau_1\pmul)(\cdot)$
and using the $\tau_1$-orthogonality
$(\xi_h, \Ponez\xi_h)_{\tau_1} = \norm{\Ponez\xi_h}_{\tau_1}^2$,
\[
\begin{aligned}
  (\bw_h, \xi_h)
  &= P_1 + \bigl((\tau_1\smul - \tau_1\pmul)\Ponez\xi_h,\ \xi_h\bigr) \\
  &\;\ge\; P_1 - \psi\,\norm{\Ponez\xi_h}_{\tau_1}\norm{\xi_h}_{\tau_1}
  \;\ge\; P_1 - \frac{2\psi}{\beta_0^2}\,(P_1 + O_1),
\end{aligned}
\]
where the last step used the nonexpansiveness of $\Ponez$ in $\norm{\cdot}_{\tau_1}$
together with \cref{H:projection}, in the form
$\norm{\xi_h}_{\tau_1}^2 \le (2/\beta_0^2)(P_1 + O_1)$. For \cref{oa:eq:T5}, since
$\delta_h$ has zero mean the shift is invisible as noted, and by the exact
$\tau_2$-orthogonal decomposition of $\delta_h$ with respect to $Q_h$,
$\norm{\delta_h}_{\tau_2}^2 = P_2 + O_2$; hence
\[
  (r_h, \delta_h)
  = P_2 + \bigl((\tau_2\smul - \tau_2\pmul)\Ptwo\delta_h,\ \delta_h\bigr)
  \;\ge\; P_2 - \psi\,(P_2 + O_2).
\]
This is the point where the mass slot enjoys the exact analogue of \cref{H:projection} with
constant one: because $\delta_h$ has zero mean, projecting onto $Q_h$ and correcting
the test function by a constant replaces the constrained projection
$\Pi_{2,0}$ of~\cite{Codina2008FiniteEA} at no cost.

\medskip\noindent\emph{Step 3: the perturbation terms.}
For \cref{oa:eq:T1}, by Cauchy--Schwarz, $|\alpha| \le \alpha_K$ elementwise,
$|\ViscProj \boldsymbol{M}| \le |\boldsymbol{M}|$ pointwise, the inverse estimate
\cref{eq:InverseEstimateFiniteOrderNorm}, \cref{oa:eq:Smoothing} and \cref{oa:eq:K1},
\begin{align*}
  \text{\cref{oa:eq:T1}}
  &\ \ge\ -\,2\sum_K \nu^{1/2}\,\bigl\lVert\alpha^{1/2}\ViscProj\nabla\bu_h\bigr\rVert_K
  \cdot \nu^{1/2}\alpha_K^{1/2}\,\frac{C_{\mathrm{inv}}}{h_K}\,(1+\psi)\,\tau_{1,K}
  \norm{\Ponez\xi_h}_K\\
  &\ \ge\ -\,2 M_1 \Bigl(\frac{D_\nu}{2}\Bigr)^{1/2} P_1^{1/2},
\end{align*}
with $M_1 \coloneqq (1+\psi)/\gamma$, whence, by Young's inequality,
\[
  \text{\cref{oa:eq:T1}} \ \ge\ -\frac{D_\nu}{4} - 2M_1^2\, P_1 .
\]
For \cref{oa:eq:T3}, using $\sigma\tau_1 \le 1$ from \cref{oa:eq:P1},
\begin{align*}
  \text{\cref{oa:eq:T3}}
  &\ \ge\ -(1+\psi)\sum_K \sigma\,\tau_{1,K}\normK{\bu_h}\,\normK{\Ponez\xi_h}\\
  &\ \ge\ -(1+\psi)\, D_\sigma^{1/2} P_1^{1/2}
  \ \ge\ -\frac{D_\sigma}{2} - \frac{(1+\psi)^2}{2}\,P_1 .
\end{align*}
For \cref{oa:eq:T4}, using $\varepsilon\tau_2 \le C_2$ from \cref{oa:eq:P1},
\begin{align*}
  \text{\cref{oa:eq:T4}}
  &\ \ge\ -(1+\psi)\sum_K \varepsilon\,\tau_{2,K}\normK{p_h}\normK{\Ptwo\delta_h}\\
  &\ \ge\ -(1+\psi)\,C_2^{1/2}\, D_\varepsilon^{1/2} P_2^{1/2}
  \ \ge\ -\frac{D_\varepsilon}{2} - \frac{(1+\psi)^2 C_2}{2}\,P_2 .
\end{align*}
For \cref{oa:eq:T6} we bound the two contributions to
$\alpha X(V^0) = \alpha\ba\cdot\nabla\bw_h + \alpha\nabla r_h$ elementwise. By the
inverse estimate, \cref{oa:eq:Smoothing} and \cref{oa:eq:K2},
\[
\begin{aligned}
  \tau_{1,K}^{1/2}\normK{\alpha\,\ba\cdot\nabla\bw_h}
  &\le \tau_{1,K}^{1/2}\,\alpha_K |\ba|_{\infty,K}\frac{C_{\mathrm{inv}}}{h_K}
  (1+\psi)\,\tau_{1,K}\normK{\Ponez\xi_h} \\
  &\le \frac{1+\psi}{\gamma}\,\tau_{1,K}^{1/2}\normK{\Ponez\xi_h},
\end{aligned}
\]
and by \cref{oa:eq:K3},
\[
  \tau_{1,K}^{1/2}\normK{\alpha\nabla r_h}
  \le \tau_{1,K}^{1/2}\,\alpha_K\,\frac{C_{\mathrm{inv}}}{h_K}(1+\psi)\,
  \tau_{2,K}\normK{\Ptwo\delta_h}
  \le \frac{1+\psi}{\gamma}\,\tau_{2,K}^{1/2}\normK{\Ptwo\delta_h}.
\]
Hence, by Cauchy--Schwarz over the elements and $ab \le (a^2+b^2)/2$,
\[
  \text{\cref{oa:eq:T6}}
  \ \ge\ -\frac{1+\psi}{\gamma}\, O_1^{1/2}\bigl( P_1^{1/2} + P_2^{1/2} \bigr)
  \ \ge\ -\frac{1+\psi}{\gamma}\Bigl( O_1 + \frac{P_1 + P_2}{2} \Bigr).
\]
For \cref{oa:eq:T7}, expanding
$\nabla\cdot(\alpha\bw_h) = \alpha\,\nabla\cdot\bw_h + \nabla\alpha\cdot\bw_h$ and
using \cref{H:porosity} in the form
$\norm{\nabla\alpha}_{L^\infty(K)} \le C_{\nabla\alpha}\,\alpha_K/h_K$,
\[
\begin{aligned}
  \tau_{2,K}^{1/2}\normK{\nabla\cdot(\alpha\bw_h)}
  &\le \Bigl(1 + \frac{C_{\nabla\alpha}}{C_{\mathrm{inv}}}\Bigr)
  \tau_{2,K}^{1/2}\,\alpha_K\,\frac{C_{\mathrm{inv}}}{h_K}\,
  (1+\psi)\,\tau_{1,K}\normK{\Ponez\xi_h} \\
  &\le \frac{M_2(1+\psi)}{\gamma}\,\tau_{1,K}^{1/2}\normK{\Ponez\xi_h},
\end{aligned}
\]
by \cref{oa:eq:K3}, with $M_2 = 1 + C_{\nabla\alpha}/C_{\mathrm{inv}}$; hence
\[
  \text{\cref{oa:eq:T7}}
  \ \ge\ -\frac{M_2(1+\psi)}{\gamma}\, O_2^{1/2} P_1^{1/2}
  \ \ge\ -\frac{M_2(1+\psi)}{2\gamma}\,\bigl( O_2 + P_1 \bigr).
\]

\medskip\noindent\emph{Step 4: collection.}
Adding \cref{oa:eq:StepDiag} and the bounds of Steps~2--3,
\begin{align*}
  B_{\mathrm{osgs}}(\ba; U_h, W)
  \ \ge\ & \tfrac34\, D_\nu + \tfrac12\, D_\sigma + \tfrac12\, D_\varepsilon\\
  &+ \Bigl( 1 - \frac{2\psi}{\beta_0^2} - \frac{1+\psi}{\gamma} \Bigr)\, O_1
  + \Bigl( 1 - \psi - \frac{M_2(1+\psi)}{2\gamma} \Bigr)\, O_2\\
  &+ \Bigl( 1 - \frac{2\psi}{\beta_0^2} - 2M_1^2 - \frac{(1+\psi)^2}{2}
  - \frac{1+\psi}{2\gamma} - \frac{M_2(1+\psi)}{2\gamma} \Bigr)\, P_1\\
  &+ \Bigl( 1 - \psi - \frac{(1+\psi)^2 C_2}{2} - \frac{1+\psi}{2\gamma} \Bigr)\, P_2 .
\end{align*}
With $\gamma \ge \gamma_0 = 10(1+M_2)$ and
$\psi \le \psi_0 = \min\{1/8, \beta_0^2/16\}$ (so that $2\psi/\beta_0^2 \le 1/8$ and
$1+\psi \le 9/8$), a direct evaluation shows that every coefficient is at least
$1/8$ (for $P_1$: $1 - 0.125 - 0.007 - 0.633 - 0.029 - 0.057 \ge 0.14$; for $P_2$,
using $C_2 \le 1$: $1 - 0.125 - 0.633 - 0.029 \ge 0.21$; the others are larger).
Hence, with $C_0 = 1/8$,
\begin{equation}
  \label{oa:eq:Collected}
  B_{\mathrm{osgs}}(\ba; U_h, W)
  \ \ge\ C_0\,\bigl( D_\nu + D_\sigma + D_\varepsilon + O_1 + O_2 + P_1 + P_2 \bigr).
\end{equation}

\medskip\noindent\emph{Step 5: recovery of the triple norm and boundedness of the test
function.}
By \cref{H:projection} and the exact decomposition of the mass slot,
\[
  \tnorm{U_h}_{\mathrm{O}}^2
  = D_\nu + D_\sigma + D_\varepsilon + \norm{\xi_h}_{\tau_1}^2 + \norm{\delta_h}_{\tau_2}^2
  \le D_\nu + D_\sigma + D_\varepsilon + \frac{2}{\beta_0^2}(P_1 + O_1) + (P_2 + O_2),
\]
so \cref{oa:eq:Collected} yields
$B_{\mathrm{osgs}}(\ba; U_h, W) \ge C_0 \min\{1, \beta_0^2/2\}\,\tnorm{U_h}_{\mathrm{O}}^2$.
It remains to bound $\tnorm{W}_{\mathrm{O}} \le \tnorm{U_h}_{\mathrm{O}} + \tnorm{V^0}_{\mathrm{O}}$. Using
\cref{oa:eq:Smoothing} elementwise together with \cref{oa:eq:K1} for the viscous term,
$\sigma\tau_1 \le 1$ for the reactive term, $\varepsilon\tau_2 \le C_2$ for the
compressibility term, and the bounds of Step~3 for
$\norm{\alpha X(V^0)}_{\tau_1}$ and $\norm{\nabla\cdot(\alpha\bw_h)}_{\tau_2}$,
\begin{align*}
  \tnorm{V^0}_{\mathrm{O}}^2
  &\le \Bigl[ \frac{2(1+\psi)^2}{\gamma^2}
  + (1+\psi)^2
  + \frac{2(1+\psi)^2}{\gamma^2}
  + \frac{M_2^2(1+\psi)^2}{\gamma^2} \Bigr] P_1 \\
  &+ \Bigl[ (1+\psi)^2 C_2 + \frac{2(1+\psi)^2}{\gamma^2} \Bigr] P_2\\
  &\le C\,(P_1 + P_2),
\end{align*}
and $P_1 \le \norm{\xi_h}_{\tau_1}^2$, $P_2 \le \norm{\delta_h}_{\tau_2}^2$ give
$\tnorm{V^0}_{\mathrm{O}} \le C\,\tnorm{U_h}_{\mathrm{O}}$, hence $\tnorm{W}_{\mathrm{O}} \le C\,\tnorm{U_h}_{\mathrm{O}}$. Combining,
\[
  \sup_{V_h\in\Xhz} \frac{B_{\mathrm{osgs}}(\ba; U_h, V_h)}{\tnorm{V_h}_{\mathrm{O}}}
  \ \ge\ \frac{B_{\mathrm{osgs}}(\ba; U_h, W)}{\tnorm{W}_{\mathrm{O}}}
  \ \ge\ \beta_{\mathrm{osgs}}\, \tnorm{U_h}_{\mathrm{O}}
\]
with $\beta_{\mathrm{osgs}} = C_0\min\{1,\beta_0^2/2\}/C > 0$, which is
\cref{oa:eq:InfSup}.
\end{proof}

\begin{remark}[On the structure of the proof]
\label{oa:rem:proofstructure}
The proof is that of~\cite[Theorem~1]{Codina2008FiniteEA} with four modifications. (i)~All
inverse-estimate steps carry the porosity weight, absorbed through
\cref{oa:eq:K1}--\cref{oa:eq:K3}; only the plain inverse estimate
\cref{eq:InverseEstimateFiniteOrderNorm} is used---in contrast with the ASGS stability proof
of~the main text, no estimate for $\nabla\cdot(\alpha\nu\ViscProj\nabla\bv_h)$ is
needed, because the analyzed OSGS form contains no viscous stabilization term.
(ii)~The reactive pairing \cref{oa:eq:T3} is new; it is absorbed by the \emph{full}
Galerkin reaction term $D_\sigma$ at the price of half of $P_1$, using nothing more
than $\sigma\tau_1 \le 1$. This is where the OSGS analysis retains
$\norm{\sigma^{1/2}\bu_h}$ in the norm, where the ASGS analysis retains only the
damped $\widetilde{\sigma}_\alpha$. (iii)~The compressibility pairing \cref{oa:eq:T4} is
new and is controlled by the same condition $\varepsilon\tau_2 \le C_2$ assumed
in~the main text. (iv)~The constrained pressure projection $\Pi_{2,0}$
of~\cite{Codina2008FiniteEA} is replaced by the unconstrained projection plus a mean
correction of the test function, which is invisible in every term; this makes the
mass-slot stability constant exactly one and is the reason no analogue of
\cref{H:projection} is required for the divergence.
\end{remark}

\subsection{Convergence}
\label{oa:sec:convergence}

Throughout this section, $U = [\bu; p]$ denotes the solution of
the linearized problem (\cref{sec:CommonSetting}) under the regularity \cref{H:regularity}, and $U_h \in \Xhz$ the solution of
\cref{oa:eq:Method}. That solution exists and is unique: for $h \le h_0$ the inf--sup estimate
\cref{oa:eq:InfSup} of \cref{oa:th:stability}, on the square linear system with equal trial and
test spaces $\Xhz$, makes the system operator injective and hence invertible, so
\cref{oa:eq:Method} is well posed. We use the local interpolation{\hyp}error quantities
$\Eint{K}{\bu}$ and $\Eint{K}{p}$ of the main text (\cref{eq:InterpolationError}), and state
the convergence estimate in the porosity-weighted broken $\ell^2$ functional
\begin{equation}
  \label{oa:eq:ErrorFunctionL2}
  \Psi_{\mathrm O}(h)
  \coloneqq
  \Bigl( \sum_K \bigl( c_1 + \Dah \bigr)\,\frac{\alpha_K^2}{h_K^2}
  \bigl( \tau_{2,K}\,\Eint{K}{\bu}^2 + \tau_{1,K}\,\Eint{K}{p}^2 \bigr) \Bigr)^{1/2},
\end{equation}
whose coarser $\ell^1$ majorant is the error function
\begin{equation}
  \label{oa:eq:ErrorFunction}
  \mathcal{E}(h)
  \coloneqq
  \sum_K \frac{\alpha_K}{h_K}\,\bigl(1 + \Dah\bigr)^{1/2}
  \Bigl( \tau_{2,K}^{1/2}\,\Eint{K}{\bu} + \tau_{1,K}^{1/2}\,\Eint{K}{p} \Bigr),
\end{equation}
i.e., the error function of the ASGS convergence theorem of~the main text with the
local factor $(1+\Dah)^{1/2}$ inserted. The two are related by
$\Psi_{\mathrm O}(h) \le c_1^{1/2}\,\mathcal{E}(h)$ (the discrete $\ell^2\subset\ell^1$
inequality together with $c_1+\Dah \le c_1(1+\Dah)$). We state the theorem in
$\Psi_{\mathrm O}$: the $\ell^1$ passage to $\mathcal{E}$ can lose a mesh-cardinality
factor $\sim h^{-d/2}$ on a quasi-uniform family, so the optimal-order corollaries
must be read off the $\ell^2$ functional.

\subsubsection{Consistency}

The analyzed method is not consistent in the classical variational sense
(cf.~\cite[Eq.~(44)]{Codina2008FiniteEA}); its consistency error is the value of the
stabilization terms at the exact solution, and it is of the order of the error
function:

\begin{lemma}[Consistency error]
\label{oa:lem:consistency}
Under \cref{H:mesh}--\cref{H:regularity}, \cref{H:design}--\cref{H:advectionsmooth}, for all $V_h \in \Xhz$,
\begin{equation}
  \label{oa:eq:ConsistencyId}
  B_{\mathrm{osgs}}(\ba;\, U - U_h,\, V_h)
  = S_1(U, V_h) + S_2(U, V_h),
\end{equation}
and
\begin{equation}
  \label{oa:eq:ConsistencyBound}
  \bigl| B_{\mathrm{osgs}}(\ba;\, U - U_h,\, V_h) \bigr|
  \;\le\; C\,\Psi_{\mathrm O}(h)\,\tnorm{V_h}_{\mathrm{O}}.
\end{equation}
\end{lemma}

\begin{proof}
Since $U$ solves the linearized problem of \cref{sec:CommonSetting} with the stated regularity, testing the
strong equations against $V_h \in \Xhz$ and integrating the viscous term by parts
gives $B(\ba; U, V_h) = L(V_h)$; adding $S_1(U,V_h) + S_2(U,V_h)$ to both sides and
subtracting \cref{oa:eq:Method} yields \cref{oa:eq:ConsistencyId}.

For the momentum part, split $\alpha X(U) = \alpha(\ba\cdot\nabla\bu) + \alpha\nabla p$
and use, respectively, Cauchy--Schwarz in $(\cdot,\cdot)_{\tau_1}$, linearity of
$\Poneo$ and \cref{oa:lem:bestapprox} with $z = \ba\cdot\nabla\bu$, $r = k_u$ and with
$z = \nabla p$, $r = k_p$:
\[
  |S_1(U, V_h)|
  \le \Bigl( \norm{\Poneo[\alpha\,\ba\cdot\nabla\bu]}_{\tau_1}
  + \norm{\Poneo[\alpha\nabla p]}_{\tau_1} \Bigr)\,\norm{\alpha X(V_h)}_{\tau_1}.
\]
By \cref{H:advectionsmooth} and the Leibniz rule,
$\norm{\ba\cdot\nabla\bu}_{H^{k_u}(K)} \le C(1+C_D)\,|\ba|_{\infty,K}\,
\norm{\bu}_{H^{k_u+1}(K)}$, so \cref{oa:eq:WeightedApprox} and the second inequality of
\cref{oa:eq:P4} give
\[
\begin{aligned}
  \norm{\Poneo[\alpha\,\ba\cdot\nabla\bu]}_{\tau_1}
  &\le C \Bigl( \sum_K \tau_{1,K}\,\alpha_K^2\,|\ba|_{\infty,K}^2\, h_K^{2k_u}
  \norm{\bu}_{H^{k_u+1}(K)}^2 \Bigr)^{1/2} \\
  &\le C \Bigl( \sum_K \Bigl(\frac{\alpha_K}{h_K}\Bigr)^{2} \tau_{2,K}\,
  \Eint{K}{\bu}^2 \Bigr)^{1/2},
\end{aligned}
\]
while $\norm{\nabla p}_{H^{k_p}(K)} \le \norm{p}_{H^{k_p+1}(K)}$ and
\cref{oa:eq:WeightedApprox} give directly
\[
\begin{aligned}
  \norm{\Poneo[\alpha\nabla p]}_{\tau_1}
  &\le C \Bigl( \sum_K \tau_{1,K}\,\alpha_K^2\, h_K^{2k_p}
  \norm{p}_{H^{k_p+1}(K)}^2 \Bigr)^{1/2} \\
  &= C \Bigl( \sum_K \Bigl(\frac{\alpha_K}{h_K}\Bigr)^{2} \tau_{1,K}\,
  \Eint{K}{p}^2 \Bigr)^{1/2}.
\end{aligned}
\]
Both right-hand sides are bounded by $C\,\Psi_{\mathrm O}(h)$ (with room to spare, since
$c_1+\Dah \ge 1$ and no Damk\"ohler factor arises here), and
$\norm{\alpha X(V_h)}_{\tau_1} \le \tnorm{V_h}_{\mathrm{O}}$.

For the mass part, the exact mass equation \cref{eq:StrongMassEquation} gives
$\nabla\cdot(\alpha\bu) = -\varepsilon p$, hence, by \cref{oa:lem:annihilation}-style
linearity,
\[
  S_2(U, V_h) = -\varepsilon\,\bigl( \Ptwoo p,\ \nabla\cdot(\alpha\bv_h)
  \bigr)_{\tau_2},
\]
which vanishes when $\varepsilon = 0$. For $\varepsilon > 0$, by Cauchy--Schwarz,
\cref{oa:lem:bestapprox} (with $\alpha$ replaced by $1$, $z = p$, $r = k_p+1$),
$\varepsilon\tau_{2} \le C_2 \le 1$ and $\varepsilon \le C_2\tau_2^{-1}$ from
\cref{oa:eq:P1}, and finally \cref{oa:eq:P5},
\begin{align*}
  |S_2(U, V_h)|
  &\le \varepsilon \norm{\Ptwoo p}_{\tau_2}\,\norm{\nabla\cdot(\alpha\bv_h)}_{\tau_2}
  \le C \Bigl( \sum_K \varepsilon^2 \tau_{2,K}\, h_K^{2(k_p+1)}
  \norm{p}_{H^{k_p+1}(K)}^2 \Bigr)^{1/2} \tnorm{V_h}_{\mathrm{O}}\\
  &\le C\,C_2 \Bigl( \sum_K \tau_{2,K}^{-1}\, \Eint{K}{p}^2 \Bigr)^{1/2}
  \tnorm{V_h}_{\mathrm{O}}\\
  &\le C \Bigl( \sum_K \bigl(c_1 + \Dah\bigr)
  \Bigl(\frac{\alpha_K}{h_K}\Bigr)^{2}\tau_{1,K}\,\Eint{K}{p}^2 \Bigr)^{1/2}
  \tnorm{V_h}_{\mathrm{O}}
\end{align*}
(in the last step, $\varepsilon\tau_2 \le C_2 \le 1$ was applied once more, followed
by \cref{oa:eq:P5}),
which is exactly $C\,\Psi_{\mathrm O}(h)\,\tnorm{V_h}_{\mathrm{O}}$. Each of the two broken $\ell^2$
right-hand sides is thus dominated by $\Psi_{\mathrm O}(h)$, so summing the two parts
proves \cref{oa:eq:ConsistencyBound} in the $\ell^2$ functional (with no $\ell^2\subset\ell^1$
enlargement).
\end{proof}

\subsubsection{Interpolation estimates}

Let $\hat U_h = [\hat\bu_h; \hat p_h]$ where $\hat\bu_h \in V_{h0}$ is the nodal
interpolant of $\bu$ (well defined since $\bu \in H^{k_u+1} \subset
C^0(\overline\Omega)$ for $d\le 3$, and boundary-value preserving) and
$\hat p_h \in Q_{h0}$ is the nodal interpolant of $p$, shifted by its mean when
$Q_{h0}$ consists of zero-mean functions. Write
$\hat E \coloneqq U - \hat U_h = [\hat{\boldsymbol e}_u; \hat e_p]$.

\begin{remark}[The pressure mean shift]
\label{oa:rem:meanshift}
Since $\int_\Omega p = 0$, the shift constant is bounded by the raw interpolation
error, $|\Omega|^{1/2}|m| = |\int_\Omega (\hat p_h - p)|\,|\Omega|^{-1/2}
\le \norm{p - \hat p_h}$, so all interpolation estimates for $\hat e_p$ hold with at
most doubled constants. Moreover the constant part of $\hat e_p$ is invisible in the
pairings below: against $\nabla\cdot(\alpha\bv_h)$ it integrates to zero by the
divergence theorem, against $\varepsilon q_h$ it vanishes because $q_h$ has zero mean
(or no shift was performed), and it has no gradient. We therefore do not distinguish
the shifted interpolant notationally.
\end{remark}

\begin{lemma}[Interpolation continuity]
\label{oa:lem:interpolation}
Under \cref{H:mesh}--\cref{H:regularity}, \cref{H:design}--\cref{H:advectionsmooth}, for all $V_h \in \Xhz$,
\begin{equation}
  \label{oa:eq:InterpolationBound}
  \bigl| B_{\mathrm{osgs}}(\ba;\, \hat E,\, V_h) \bigr|
  \;\le\; C\,\Psi_{\mathrm O}(h)\,\tnorm{V_h}_{\mathrm{O}}.
\end{equation}
\end{lemma}

\begin{proof}
Expand $B_{\mathrm{osgs}}(\ba; \hat E, V_h)$ into the Galerkin terms and the
stabilization terms and estimate each. Below, all interpolation estimates are
\cref{oa:eq:NodalInterp} with $n = k_u+1$ or $n = k_p+1$ and $m \in \{0,1\}$.

\smallskip\noindent
\emph{(I1) Viscous term.} Using \cref{oa:eq:P4} ($\nu \le \alpha_K\tau_{2,K}$),
\begin{align*}
  2\nu(\alpha\ViscProj\nabla\hat{\boldsymbol e}_u, \ViscProj\nabla\bv_h)
  &\le \sqrt{2}\,\tnorm{V_h}_{\mathrm{O}}\Bigl( \sum_K 2\nu\alpha_K\,
  \norm{\nabla\hat{\boldsymbol e}_u}_K^2 \Bigr)^{1/2}\\
  &\le C\,\tnorm{V_h}_{\mathrm{O}}\Bigl( \sum_K
  \Bigl(\frac{\alpha_K}{h_K}\Bigr)^2\tau_{2,K}\,\Eint{K}{\bu}^2 \Bigr)^{1/2}.
\end{align*}

\smallskip\noindent
\emph{(I2) Convective and mass--divergence terms, jointly.} Using
$\nabla\cdot(\alpha\ba) = 0$, $\hat{\boldsymbol e}_u \in H^1_0(\Omega)^d$
(nodal interpolation preserves the boundary values of $\bu$) and the continuity of the
pressure space, global integration by parts gives
\[
  (\bv_h, \alpha\ba\cdot\nabla\hat{\boldsymbol e}_u)
  + (q_h, \nabla\cdot(\alpha\hat{\boldsymbol e}_u))
  = -\bigl( \hat{\boldsymbol e}_u,\ \alpha\bigl(\ba\cdot\nabla\bv_h + \nabla q_h\bigr)
  \bigr)
  = -\bigl( \hat{\boldsymbol e}_u,\ \alpha X(V_h) \bigr).
\]
Hence, splitting $\tau_{1,K}^{-1/2} \le (\alpha_K\tns^{-1})^{1/2} + \sigma^{1/2}$ and
using \cref{oa:eq:P3} for the first part and \cref{oa:eq:P6} for the second,
\[
\begin{aligned}
  \bigl|\bigl(\hat{\boldsymbol e}_u, \alpha X(V_h)\bigr)\bigr|
  &\le \tnorm{V_h}_{\mathrm{O}} \Bigl( \sum_K \tau_{1,K}^{-1} \norm{\hat{\boldsymbol e}_u}_K^2
  \Bigr)^{1/2} \\
  &\le C\,\tnorm{V_h}_{\mathrm{O}}\Bigl( \sum_K \bigl(c_1 + \Dah\bigr)
  \Bigl(\frac{\alpha_K}{h_K}\Bigr)^2 \tau_{2,K}\,\Eint{K}{\bu}^2 \Bigr)^{1/2}.
\end{aligned}
\]
This is the first of the two places where the Damk\"ohler factor enters.

\smallskip\noindent
\emph{(I3) Reactive term.} Using the full reactive control in the norm
\cref{oa:eq:TripleNorm} and \cref{oa:eq:P6},
\[
\begin{aligned}
  \sigma(\hat{\boldsymbol e}_u, \bv_h)
  &\le \norm{\sigma^{1/2}\bv_h}\Bigl( \sum_K \sigma\,
  \norm{\hat{\boldsymbol e}_u}_K^2 \Bigr)^{1/2} \\
  &\le C\,\tnorm{V_h}_{\mathrm{O}}\Bigl( \sum_K \Dah
  \Bigl(\frac{\alpha_K}{h_K}\Bigr)^2\tau_{2,K}\,\Eint{K}{\bu}^2 \Bigr)^{1/2}.
\end{aligned}
\]

\smallskip\noindent
\emph{(I4) Compressibility term.} Using $\varepsilon \le C_2\tau_2^{-1}$ and
\cref{oa:eq:P5},
\[
\begin{aligned}
  \varepsilon(\hat e_p, q_h)
  &\le \varepsilon^{1/2}\norm{\varepsilon^{1/2}q_h}
  \Bigl( \sum_K \norm{\hat e_p}_K^2 \Bigr)^{1/2} \\
  &\le C\,\tnorm{V_h}_{\mathrm{O}}\Bigl( \sum_K \bigl(c_1+\Dah\bigr)
  \Bigl(\frac{\alpha_K}{h_K}\Bigr)^2\tau_{1,K}\,\Eint{K}{p}^2 \Bigr)^{1/2}.
\end{aligned}
\]

\smallskip\noindent
\emph{(I5) Pressure-gradient term.} Global integration by parts
($\bv_h \in V_{h0}$) and \cref{oa:eq:P5} give
\[
\begin{aligned}
  (\bv_h, \alpha\nabla\hat e_p)
  &= -(\hat e_p, \nabla\cdot(\alpha\bv_h)) \\
  &\le \norm{\nabla\cdot(\alpha\bv_h)}_{\tau_2}
  \Bigl( \sum_K \tau_{2,K}^{-1}\norm{\hat e_p}_K^2 \Bigr)^{1/2} \\
  &\le C\,\tnorm{V_h}_{\mathrm{O}}\Bigl( \sum_K \bigl(c_1+\Dah\bigr)
  \Bigl(\frac{\alpha_K}{h_K}\Bigr)^2\tau_{1,K}\,\Eint{K}{p}^2 \Bigr)^{1/2},
\end{aligned}
\]
the second place where the Damk\"ohler factor enters (through
$\tau_2^{-1/2} \le (c_1+\Dah)^{1/2}(\alpha_K/h_K)\tau_1^{1/2}$).

\smallskip\noindent
\emph{(I6) Momentum stabilization.} Since fluctuation operators are nonexpansive,
\[
  |S_1(\hat E, V_h)|
  \le \bigl( \norm{\alpha\,\ba\cdot\nabla\hat{\boldsymbol e}_u}_{\tau_1}
  + \norm{\alpha\nabla\hat e_p}_{\tau_1} \bigr)\,\tnorm{V_h}_{\mathrm{O}}.
\]
Elementwise, by \cref{oa:eq:P4},
\[
\begin{aligned}
  \tau_{1,K}^{1/2}\,\alpha_K\,|\ba|_{\infty,K}\norm{\nabla\hat{\boldsymbol e}_u}_K
  &\le C\,\frac{\alpha_K}{h_K}\,\tau_{2,K}^{1/2}\,\Eint{K}{\bu},
  \qquad \\
  \tau_{1,K}^{1/2}\,\alpha_K\norm{\nabla\hat e_p}_K
  &\le C\,\frac{\alpha_K}{h_K}\,\tau_{1,K}^{1/2}\,\Eint{K}{p},
\end{aligned}
\]
the second bound holding exactly; no Damk\"ohler factor arises.

\smallskip\noindent
\emph{(I7) Mass stabilization.} Likewise,
$|S_2(\hat E, V_h)| \le
\norm{\nabla\cdot(\alpha\hat{\boldsymbol e}_u)}_{\tau_2}\tnorm{V_h}_{\mathrm{O}}$, and elementwise,
using \cref{H:porosity} in the form $\norm{\nabla\alpha}_{L^\infty(K)} \le C_{\nabla\alpha}\alpha_K/h_K$,
\[
  \tau_{2,K}^{1/2}\norm{\nabla\cdot(\alpha\hat{\boldsymbol e}_u)}_K
  \le \tau_{2,K}^{1/2}\bigl( \alpha_K \norm{\nabla\hat{\boldsymbol e}_u}_K
  + C_{\nabla\alpha}\tfrac{\alpha_K}{h_K}\norm{\hat{\boldsymbol e}_u}_K \bigr)
  \le C(1+C_{\nabla\alpha})\,\frac{\alpha_K}{h_K}\,\tau_{2,K}^{1/2}\,\Eint{K}{\bu}.
\]

Collecting (I1)--(I7), each square-summed right-hand side is dominated by
$\Psi_{\mathrm O}(h)$---the momentum, reaction and viscous groups through
$c_1+\Dah \ge 1$, the convective, compressibility and pressure-gradient groups through their explicit
$(c_1+\Dah)$ factor---and, absorbing $c_1$-dependent constants into $C$, we obtain
\cref{oa:eq:InterpolationBound} directly in the $\ell^2$ functional.
\end{proof}

\begin{lemma}[Size of the interpolation error in the triple norm]
\label{oa:lem:interpsize}
Under \cref{H:mesh}--\cref{H:regularity}, \cref{H:design}--\cref{H:advectionsmooth},
$\tnorm{\hat E}_{\mathrm{O}} \le C\,\Psi_{\mathrm O}(h)$.
\end{lemma}

\begin{proof}
Estimate the five terms of \cref{oa:eq:TripleNorm} at $V = \hat E$: the viscous term as
in (I1); $\norm{\sigma^{1/2}\hat{\boldsymbol e}_u}$ as in (I3);
$\varepsilon^{1/2}\norm{\hat e_p}$ as in (I4);
$\norm{\alpha X(\hat E)}_{\tau_1}$ as in (I6); and
$\norm{\nabla\cdot(\alpha\hat{\boldsymbol e}_u)}_{\tau_2}$ as in (I7).
\end{proof}

\subsubsection{Main theorem and corollaries}

\begin{theorem}[Convergence of the OSGS method]
\label{oa:th:convergence}
Let $U$ solve the linearized problem (\cref{sec:CommonSetting}) and let
$U_h \in \Xhz$ solve \cref{oa:eq:Method}. Under \cref{H:mesh}--\cref{H:regularity}, \cref{H:design}--\cref{H:advectionsmooth} and for
$h \le h_0$ as in \cref{oa:th:stability}, there is a constant $C$, independent of $h$,
$\nu$, $\sigma$, $\varepsilon$ and of $\alpha$, $\ba$ except through the constants of
\cref{oa:sec:hypotheses}, such that
\begin{equation}
  \label{oa:eq:MainResult}
\begin{aligned}
  \tnorm{U - U_h}_{\mathrm{O}}
  &\;\le\;
  C\,\Psi_{\mathrm O}(h) \\
  &= C\Bigl( \sum_K \bigl( c_1 + \Dah \bigr)\,\frac{\alpha_K^2}{h_K^2}
  \bigl( \tau_{2,K}\,\Eint{K}{\bu}^2 + \tau_{1,K}\,\Eint{K}{p}^2 \bigr) \Bigr)^{1/2} \\
  &\;\le\; C\,\mathcal{E}(h),
\end{aligned}
\end{equation}
with $\Psi_{\mathrm O}(h)$ and its $\ell^1$ majorant $\mathcal{E}(h)$ as in
\cref{oa:eq:ErrorFunctionL2,oa:eq:ErrorFunction}.
\end{theorem}

\begin{proof}
Write $U - U_h = \hat E + \eta_h$ with $\eta_h \coloneqq \hat U_h - U_h \in \Xhz$.
By \cref{oa:th:stability} there is $V_h \in \Xhz$, $\tnorm{V_h}_{\mathrm{O}} = 1$, with
\[
  \beta_{\mathrm{osgs}}\,\tnorm{\eta_h}_{\mathrm{O}}
  \le B_{\mathrm{osgs}}(\ba; \eta_h, V_h)
  = -\,B_{\mathrm{osgs}}(\ba; \hat E, V_h)
  + B_{\mathrm{osgs}}(\ba; U - U_h, V_h)
  \le C\,\Psi_{\mathrm O}(h),
\]
by \cref{oa:lem:interpolation,oa:lem:consistency}. The triangle inequality and
\cref{oa:lem:interpsize} conclude, both in the $\ell^2$ functional $\Psi_{\mathrm O}(h)$.
\end{proof}

\begin{corollary}[Estimate in the ASGS norm]
\label{oa:cor:asgsnorm}
Under the hypotheses of \cref{oa:th:convergence}, the OSGS solution also satisfies the
estimate of the ASGS convergence theorem of~the main text up to the local factor
$(c_1+\Dah)^{1/2}$, in the same porosity-weighted $\ell^2$ form:
\[
\begin{aligned}
  \tnorm{U - U_h}_{\mathrm{S}}
  &\;\le\;
  C\,\Psi_{\mathrm O}(h) \\
  &= C\Bigl( \sum_K \bigl( c_1 + \Dah \bigr)\,\frac{\alpha_K^2}{h_K^2}
  \bigl( \tau_{2,K}\,\Eint{K}{\bu}^2 + \tau_{1,K}\,\Eint{K}{p}^2 \bigr) \Bigr)^{1/2} .
\end{aligned}
\]
\end{corollary}

\begin{proof}
Combine \cref{oa:th:convergence} with \cref{oa:lem:normcomparison}.
\end{proof}

\begin{corollary}[$\sigma$-robust $L^2$ velocity estimate]
\label{oa:cor:Ltwo}
Let $\sigma > 0$. Under the hypotheses of \cref{oa:th:convergence},
\begin{equation}
  \label{oa:eq:LtwoVelocity}
  \norm{\bu - \bu_h}
  \;\le\;
  C \Bigl[ \sum_K \Bigl( 1 + \frac{c_1}{\Dah} \Bigr)
  \Bigl(
  \bigl( 1 + \tfrac{c_2}{c_1}\Reh \bigr)\, \Eint{K}{\bu}^2
  + \frac{\alpha_K}{\sigma\nu}\, \Eint{K}{p}^2
  \Bigr) \Bigr]^{1/2}.
\end{equation}
The factor is written $1+c_1/\Dah$, not $1+\Dah^{-1}$: dividing $\Psi_{\mathrm O}$ by
$\sigma^{1/2}$ carries the full $c_1+\Dah$, and absorbing $c_1$ into that factor would hide
its dependence. In particular, wherever $\Dah \gtrsim c_1$ (the reaction-dominated
pre-asymptotic regime), the local factor $1+c_1/\Dah$ is $O(1)$ and, on a quasi-uniform
mesh, the $L^2$ velocity error is of full interpolation order $h^{k_u+1}$ with a constant
uniform in $\sigma$; the pressure contribution moreover \emph{decreases} as
$(\sigma\nu)^{-1/2}$ as the reaction strengthens.
\end{corollary}

\begin{proof}
Since the norm \cref{oa:eq:TripleNorm} controls $\norm{\sigma^{1/2}(\bu-\bu_h)}$,
$\norm{\bu - \bu_h} \le \sigma^{-1/2}\tnorm{U - U_h}_{\mathrm{O}} \le
C\,\sigma^{-1/2}\Psi_{\mathrm O}(h)$ by \cref{oa:th:convergence}. In each summand of
$\sigma^{-1}\Psi_{\mathrm O}(h)^2$, use $\sigma^{-1} = h_K^2/(\Dah\,\alpha_K\nu)$ to write
\[
  \sigma^{-1}\,(c_1+\Dah)\,\frac{\alpha_K^2}{h_K^2}
  = \Bigl(1 + \frac{c_1}{\Dah}\Bigr)\,\frac{\alpha_K}{\nu};
\]
then \cref{oa:eq:Tau2Expanded} gives $(\alpha_K/\nu)\,\tau_{2,K} = 1 + (c_2/c_1)\Reh$ for the
velocity term, while $\tau_{1,K} \le 1/\sigma$ from \cref{oa:eq:P1} gives
$(\alpha_K/\nu)\,\tau_{1,K} \le \alpha_K/(\sigma\nu)$ for the pressure term. Substituting
these into $\sigma^{-1/2}\Psi_{\mathrm O}(h)$ yields \cref{oa:eq:LtwoVelocity}. \qedhere
\end{proof}

\subsection{Discussion}
\label{oa:sec:discussion}

\begin{remark}[The mechanism: $\sigma$ versus $\widetilde{\sigma}_\alpha$]
\label{oa:rem:mechanism}
The single structural difference between the two analyses is the fate of the reactive
term. In the ASGS quadratic form the stabilization contains the reactive residual, and
the diagonal reactive contribution is
$\sigma\norm{\bu_h}_K^2 - \sigma^2\tau_{1}\norm{\bu_h}_K^2
= \sigma(1-\sigma\tau_1)\norm{\bu_h}_K^2$, with
$\sigma(1-\sigma\tau_1) = \sigma\,\alpha_K\tns^{-1}\tau_1 \eqsim
\widetilde{\sigma}_\alpha$: the subscales \emph{renormalize} the reaction down to
$\widetilde{\sigma}_\alpha \le \alpha_K\tns^{-1} = c_1\alpha_K^2\tau_2/h_K^2$, a
quantity always commensurate with the divergence-stabilization scale, so all reactive
pairings in the ASGS convergence proof can be traded against
$(\alpha_K/h_K)\tau_2^{1/2}$-weighted quantities with no Damk\"ohler penalty. In the
OSGS variant the orthogonal projection deletes the reactive residual identically
(\cref{oa:lem:annihilation}), so no renormalization occurs: the diagonal keeps the full
$\norm{\sigma^{1/2}\bu_h}^2$ (a stronger left-hand side, cf.\
\cref{oa:lem:normcomparison}), but the reactive and pressure interpolation pairings must
be bounded head-on, and the best conversions available to the present proof---\cref{oa:eq:P5}
and \cref{oa:eq:P6}---cost factors $(c_1+\Dah)^{1/2}$ and $\Dah^{1/2}\le(c_1+\Dah)^{1/2}$, respectively. The Damk\"ohler factor in
\cref{oa:eq:MainResult} is thus the factor these conversions introduce when the reactive
subscale is deleted, while the strengthened norm (in particular \cref{oa:cor:Ltwo}) is the
reaction control retained in return. Both are properties of the present (one-sided)
estimate, not established as sharp or necessary.
\end{remark}

\begin{remark}[Consistency with the numerical campaign of~the main text]
\label{oa:rem:numerics}
The reaction-dominated experiments of~the main text (manufactured
solutions, $\mathrm{Da} = 10^6$) report two facts about the OSGS variant relative to the ASGS
one: (i)~the $L^2$ velocity errors of the two methods are indistinguishable at all
resolutions, and (ii)~the $H^1$-type velocity error of the OSGS method exceeds the
ASGS one by a factor of about $3.5$ on the finest $P_1$ meshes---the largest velocity
discrepancy of the whole campaign---with the gap closing under refinement at an
accelerating relative rate. Both observations are compatible in order of magnitude with
the present, one-sided bounds. Fact (i) is \cref{oa:cor:Ltwo}: in the regime $\Dah \gg 1$ the
OSGS $L^2$ velocity estimate is of full order with a $\sigma$-uniform constant, so no excess
is predicted in $L^2$ (a one-sided bound cannot, of course, forbid one). Fact (ii) is consistent with the
factor $(1+\Dah)^{1/2}$ in \cref{oa:eq:MainResult}: on the finest meshes of those
experiments the global mesh Damk\"ohler number is of order ten (its precise value
depends on the element-length convention entering $\tau$), so the bound permits an
excess of about $(1+\Dah)^{1/2} \approx 3$; on the interior plateau of the porosity
profile the elementwise values are larger and permit factors of $4$--$5$. Moreover,
$\Dah^{1/2} = (\sigma/(\alpha_K\nu))^{1/2}\, h_K$ decays \emph{linearly} in $h$ while
the bound itself decays like $h^{k_u}$, which reproduces the observed accelerating
relative decay of the gap. Two caveats are mandatory. First, upper bounds are
one-sided: the analysis shows that the observed excess is \emph{permitted} and
supplies a mechanism for it (\cref{oa:rem:mechanism}); it does not prove that the excess
must occur, and the absence of the factor from the ASGS bound does not by itself prove
ASGS superiority. Second, the near-coincidence of the permitted and observed factors
at the finest resolution should be read as an order-of-magnitude consistency check,
not a validated prediction, since the constants in \cref{oa:eq:MainResult} are not
tracked sharply.
\end{remark}

\begin{remark}[Hypotheses dispensed with]
\label{oa:rem:dispensed}
Two technical hypotheses of the ASGS continuity analysis of~the main text are not
needed here: the interior-face jump condition on $\tau_1$ (required there when
$\sigma > 0$, following Eq.~(39) of~\cite{codina2001stabilized}), and the enlarged
inverse-estimate constant for quantities of the form
$\nabla\cdot(\alpha\nu\ViscProj\nabla\bv_h)$. The reason is structural: the analyzed
OSGS form contains no viscous or reactive stabilization terms, so the proof never
integrates a stabilization term by parts elementwise---the only integrations by parts
are global ones, in (I2) and (I5) of \cref{oa:lem:interpolation}. In exchange, the OSGS
analysis needs the projection stability condition \cref{H:projection}, which has no ASGS
counterpart. The two hypothesis sets are thus not nested; they reflect the different
routes (coercivity for ASGS, inf--sup for OSGS).
\end{remark}

\begin{remark}[A Method II variant]
\label{oa:rem:methodII}
As in~\cite[\S 2.2]{Codina2008FiniteEA}, one may control the convective and pressure-gradient
fluctuations separately, replacing $S_1$ by
\[
  S_1^{\mathrm{II}}(U,V) \coloneqq
  \bigl( \Poneo[\alpha\,\ba\cdot\nabla\bu],\ \alpha\,\ba\cdot\nabla\bv \bigr)_{\tau_1}
  + \bigl( \Poneo[\alpha\nabla p],\ \alpha\nabla q \bigr)_{\tau_1}.
\]
The diagonal then controls the two fluctuations individually, and the same arguments
(with the elementary inequality $a^2 + b^2 \ge \tfrac13(a^2+b^2) + \tfrac13(a+b)^2$
used in~\cite[Theorem~3]{Codina2008FiniteEA}) are expected to give stability in the correspondingly
finer norm and a convergence estimate of the same form as \cref{oa:eq:MainResult}. We state
this only as a prospective observation---the modified proof is not written out here---having
analyzed Method~I because it is the variant implemented in~the main text (a single
projected quantity per equation).
\end{remark}

\begin{remark}[Robustness regimes and the screening length]
\label{oa:rem:robustness}
Repeating the regime normalizations of the robustness section of~the main text
with \cref{oa:eq:MainResult} in place of the ASGS estimate simply multiplies each regime
estimate by $(1+\Dah)^{1/2}$. In the viscous- and convection-dominated regimes
$\Dah \le 1$ on any reasonable mesh and the two variants obey identical bounds up to
constants; asymptotically, $\Dah \propto h_K^2$, so the OSGS bound exceeds the ASGS
one by $1 + O(h_K^2)$ only. The factor matters solely in the pre-asymptotic
reaction-dominated regime, and it has a clean physical reading: introducing the
porous screening length $\ell_\sigma \coloneqq (\alpha_K\nu/\sigma)^{1/2}$ (the decay
length of velocity perturbations in a Brinkman medium), one has
$\Dah = (h_K/\ell_\sigma)^2$, so the penalty persists exactly while the mesh
under-resolves $\ell_\sigma$ and disappears once $h_K \lesssim \ell_\sigma$. In the
$\mathrm{Da} = 10^6$ experiments of~the main text, $\ell_\sigma \sim 10^{-3}L$,
below the finest resolutions employed, which is why the excess is still visible there.
\end{remark}

\begin{remark}[What this analysis does not show]
\label{oa:rem:open}
(i)~The numerical campaign of~the main text observes, for $k \ge 2$, an $L^2$
pressure convergence order for the OSGS variant one unit \emph{higher} than for the
ASGS variant. Nothing in the present analysis (nor in~\cite{Codina2008FiniteEA}, whose norms
contain no $L^2$ pressure term except in the viscous-dominated variant with
Reynolds-degrading constants) explains this; a duality argument exploiting the
orthogonality of the consistency error to the finite element space would be the
natural route, and remains open. (ii)~The analysis takes the projection equation
solved exactly; the lagged projection of the implemented staggered iteration, and the
nonlinear problem, are not covered. (iii)~The estimate \cref{oa:eq:MainResult} is an
upper bound; no lower bound is claimed, so the comparison with the ASGS estimate in
\cref{oa:cor:asgsnorm} orders the \emph{bounds}, not the methods. (iv)~For $\alpha \equiv 1$
and $\sigma = \varepsilon = 0$ the present theorem reduces to
\cite[Theorem~2]{Codina2008FiniteEA} with the parameter dictionary $c_3 = 1$,
$c_4 = c_2/c_1$, which serves as a consistency check of the constants.
\end{remark}


